% 繁體中文版的書籍排版設定；版型與共用模板保持一致。
\makeatletter
\@ifundefined{Shaded}{\newenvironment{Shaded}{}{}}{}
\makeatother
\NewCommandCopy\InfraOriginalSetMonoFont\setmonofont
\RenewDocumentCommand{\setmonofont}{O{} m}{%
  \IfFontExistsTF{#2}{\InfraOriginalSetMonoFont[#1]{#2}}{%
    \InfraOriginalSetMonoFont[#1]{DejaVu Sans Mono}}}
\input{template/agent-book-preamble.tex}
\RenewCommandCopy\setmonofont\InfraOriginalSetMonoFont
\usepackage{longtable,booktabs,array,calc}
\usepackage{xurl}
\urlstyle{same}
\setlength{\emergencystretch}{3em}
\providecommand{\BookEdition}{v1.0}
\setkeys{Gin}{width=!,height=!,keepaspectratio}
\newcommand{\infragraphic}[1]{\includegraphics[width=.68\linewidth,height=.38\textheight,keepaspectratio]{#1}}
\renewenvironment{figure}[1][!htbp]{\origfigure[#1]}{\origendfigure}
\setcounter{topnumber}{3}
\setcounter{bottomnumber}{2}
\setcounter{totalnumber}{5}
\renewcommand{\topfraction}{.85}
\renewcommand{\bottomfraction}{.75}
\renewcommand{\textfraction}{.12}
\renewcommand{\floatpagefraction}{.85}
\setlength{\floatsep}{9pt plus 2pt minus 2pt}
\setlength{\textfloatsep}{11pt plus 3pt minus 3pt}
\setlength{\intextsep}{9pt plus 2pt minus 2pt}
\captionsetup{font={small,sf,cjkmatch},justification=justified,singlelinecheck=false,skip=4pt}
\newenvironment{infratable}{\begingroup\small
  \setlength{\tabcolsep}{2.5pt}
  \setlength{\LTleft}{0pt}\setlength{\LTright}{0pt}
  \setlength{\LTpre}{6pt}\setlength{\LTpost}{6pt}
  \renewcommand{\arraystretch}{1.08}
}{\endgroup}
\allowdisplaybreaks[1]
\AtBeginDocument{%
  \setlength{\abovedisplayskip}{8pt plus 2pt minus 3pt}%
  \setlength{\belowdisplayskip}{8pt plus 2pt minus 3pt}%
  \setlength{\abovedisplayshortskip}{4pt plus 2pt}%
  \setlength{\belowdisplayshortskip}{5pt plus 2pt minus 2pt}%
}
\raggedbottom
\newunicodechar{→}{\ensuremath{\rightarrow}}
\newunicodechar{←}{\ensuremath{\leftarrow}}
\newunicodechar{↔}{\ensuremath{\leftrightarrow}}
\newunicodechar{×}{\ensuremath{\times}}
\newunicodechar{÷}{\ensuremath{\div}}
\newunicodechar{≤}{\ensuremath{\leq}}
\newunicodechar{≥}{\ensuremath{\geq}}
\newunicodechar{−}{\ensuremath{-}}
\newunicodechar{∈}{\ensuremath{\in}}
\newunicodechar{∝}{\ensuremath{\propto}}
\newunicodechar{∑}{\ensuremath{\sum}}
\newunicodechar{∞}{\ensuremath{\infty}}
\newunicodechar{²}{\ensuremath{^2}}
\newunicodechar{³}{\ensuremath{^3}}
\newunicodechar{ₖ}{\ensuremath{_k}}
\newunicodechar{ᵀ}{\ensuremath{^{\mathsf T}}}

% 優先使用 macOS 的繁中字型；CI 沒有 Apple 字型時退回 CJK sans。
\IfFontExistsTF{Songti TC}{
  \setCJKmainfont[Scale=1,BoldFont=Heiti TC]{Songti TC}
}{
  \IfFontExistsTF{Noto Sans CJK TC}{
    \setCJKmainfont[Scale=1,BoldFont=Noto Sans CJK TC]{Noto Sans CJK TC}
  }{
    \setCJKmainfont[Scale=1,BoldFont=SourceHanSansCN-Bold.otf,BoldFeatures={Path=../manuscripts/figure_style/fonts/}]{Noto Sans CJK SC}
  }
}
\IfFontExistsTF{Heiti TC}{
  \setCJKsansfont[Scale=1,BoldFont=Heiti TC]{Heiti TC}
  \setCJKfamilyfont{capcjksf}[Scale=1,BoldFont=Heiti TC]{Heiti TC}
}{
  \setCJKsansfont[Path=../manuscripts/figure_style/fonts/,Scale=1,BoldFont=SourceHanSansCN-Bold.otf]{SourceHanSansCN-Regular.otf}
  \setCJKfamilyfont{capcjksf}[Path=../manuscripts/figure_style/fonts/,Scale=1,BoldFont=SourceHanSansCN-Bold.otf]{SourceHanSansCN-Regular.otf}
}
\DeclareCaptionFont{cjkmatch}{\CJKfamily{capcjksf}}
\newunicodechar{ℓ}{\ensuremath{\ell}}
\newunicodechar{𝜇}{\ensuremath{\mu}}
\makeatletter
\setlength{\@fptop}{0pt}
\setlength{\@fpsep}{12pt plus 2pt minus 2pt}
\setlength{\@fpbot}{0pt plus 1fil}
\makeatother
\setlength{\skip\footins}{16pt plus 2pt minus 3pt}
